\section{Mise en place Kiro dans le dépôt}

Tout ce qui suit est déjà créé dans le dépôt, en anglais pour le jury. Le dossier \file{.kiro/} est à la racine et doit rester versionné.

\subsection{Steering : les règles que l'agent applique}

{\small\setlength{\tabcolsep}{4pt}
\begin{xltabular}{\textwidth}{@{}L{3.5cm}L{2.6cm}Y@{}}
\toprule
\textbf{Fichier} & \textbf{Inclusion} & \textbf{Contenu} \\
\midrule
\file{product.md} & toujours & Vision, vocabulaire (cœur, adaptateur, profil, tuteur, provenance\dots), promesse produit, règles non négociables, cible de la démo \\
\file{tech.md} & toujours & Pile technique, variables d'environnement, commandes (marquées « prévues » tant que la tâche qui les crée n'est pas faite), notes Windows \\
\file{structure.md} & toujours & Arborescence, responsables par dossier, règles de frontière, nommage, processus de changement de contrat, branches et messages de commit \\
\file{trust-and-safety.md} & toujours & Provenance, niveaux de risque, permissions, annulation, sécurité locale, vie privée, audit \\
\file{testing.md} & toujours & Couches de tests, identifiants des propriétés (\code{CON-P}, \code{CORE-P}\dots), règles (pas de vrai LLM en test, pas de contrat simulé) \\
\file{adapter-contract.md} & fichiers \file{apps/}, \file{sdk/}, contrat, passerelle & Guide court du contrat pour écrire un adaptateur, règles propres à Blender et à Chrome \\
\file{pedagogy.md} & profils, pédagogie, vérification, orchestrateur & Stratégie de tutorat, format des profils, exercices, preuves, maîtrise \\
\file{panel-ui.md} & \file{panel/}, tableau de bord & Modes d'affichage, identité visuelle, langage de provenance, règles d'interaction, accessibilité \\
\file{demo-check.md} & manuelle (\code{/demo-check}) & Liste de contrôle avant tournage ; remplace les hooks manuels absents de Kiro IDE 1.x \\
\file{release-checklist.md} & manuelle & Liste de soumission Kiroween 2026 \\
\bottomrule
\end{xltabular}}

\subsection{Specs : six périmètres, une traçabilité complète}

Chaque spec suit le format du lauréat 2025 : exigences en EARS, design avec propriétés de correction reliées aux exigences, tâches reliées aux exigences. Nous ajoutons pour chaque tâche un \textbf{responsable} et un \textbf{jalon}, et les tâches de propriétés sont obligatoires.

{\small\setlength{\tabcolsep}{4pt}
\begin{xltabular}{\textwidth}{@{}L{4.3cm}L{1.3cm}r r r r@{}}
\toprule
\textbf{Spec} & \textbf{Pilote} & \textbf{Exig.} & \textbf{Critères} & \textbf{Propr.} & \textbf{Tâches (sous-t.)} \\
\midrule
\file{adapter-contract} & \PA & 12 & 53 & 11 & 13 (7) \\
\file{core-orchestrator} & \PA & 11 & 47 & 12 & 14 (6) \\
\file{pedagogy-verification} & \PA & 6 & 26 & 9 & 9 (4) \\
\file{tutor-panel} & \PB & 6 & 28 & 6 & 11 (4) \\
\file{blender-tutor} & \PC & 8 & 31 & 5 & 12 (4) \\
\file{web-tutor} & \PC{} (\PB) & 9 & 32 & 6 & 12 (4) \\
\midrule
\textbf{Total} & & \textbf{52} & \textbf{217} & \textbf{49} & \textbf{71 (29)} \\
\bottomrule
\end{xltabular}}

Trois tâches seulement sont optionnelles (marquées \code{*}) : un adaptateur de niveau visuel, un fournisseur Bedrock et le pont MCP de développement. Elles ne sont commencées que si le planning le permet.

\subsection{Hooks : contrôles automatiques}

Les hooks sont au format actuel \file{.kiro/hooks/*.json} (version v1), lu par Kiro IDE 1.x et Kiro CLI. Les contrôles fréquents sont des \textbf{commandes} déterministes qui appellent \file{scripts/kiro/hook.py} (Python standard, testé) : elles ne consomment pas de crédits et ne font rien tant que la tâche qui crée leur cible n'est pas terminée. Les prompts d'agent sont réservés à la revue de sécurité et à l'intégration d'une nouvelle action.

{\small\setlength{\tabcolsep}{4pt}
\begin{xltabular}{\textwidth}{@{}L{4.2cm}L{2.4cm}L{1.4cm}Y@{}}
\toprule
\textbf{Hook} & \textbf{Déclencheur} & \textbf{Type} & \textbf{Effet} \\
\midrule
\code{core-boundary-guard} & sauvegarde \file{core/*.py} & commande & Refuse les imports d'hôte ou d'adaptateur et le vocabulaire d'application dans le cœur \\
\code{trust-review} & sauvegarde \file{trust/}, \file{actions/}, manifestes & agent & Revue contre les règles de confiance : risque, annulation, pas d'évaluation de code, pas de fuite de secret, permissions minimales \\
\code{task-preflight} & avant une tâche & agent & Rappelle exigences, responsable, dépendances et processus de contrat \\
\code{contract-schema-sync} & sauvegarde du contrat & commande & Régénère JSON Schema et types TypeScript \\
\code{adapter-contract-offline} & sauvegarde d'un adaptateur & commande & Kit de contrat hors ligne pour le tuteur touché \\
\code{new-action-integration} & création d'un fichier d'action & agent & Câble la nouvelle action : registre, annulation, fixture du kit, README (équivalent de leurs hooks de gestes) \\
\code{format-python}, \code{format-typescript} & sauvegarde & commande & ruff et prettier sur les fichiers concernés \\
\code{profile-validate} & sauvegarde d'un profil & commande & \code{wysp validate-profile} \\
\code{task-tests} & après une tâche & commande & Tests des paquets modifiés (pytest, Vitest) \\
\bottomrule
\end{xltabular}}

\begin{notebox}
\small Les hooks de fichiers ne se déclenchent que pour les modifications faites par l'agent, pas pour une sauvegarde manuelle dans l'éditeur. La CI exécute donc les mêmes contrôles (\code{python scripts/kiro/hook.py boundaries}, tests, export des schémas) pour couvrir le travail fait à la main.
\end{notebox}

\subsection{Agents spécialisés et MCP}

{\small\setlength{\tabcolsep}{4pt}
\begin{xltabular}{\textwidth}{@{}L{3.4cm}Y L{3.6cm}@{}}
\toprule
\textbf{Agent} & \textbf{Rôle} & \textbf{Outils} \\
\midrule
\code{blender-adapter} & Implémente et teste l'adaptateur Blender et le SDK Python, avec les règles de fil principal, d'annulation et d'événements & lecture, écriture, shell, web, Context7 \\
\code{browser-adapter} & Implémente et teste l'extension Chrome et le SDK TypeScript, avec les règles MV3, de localisateurs et de vie privée & lecture, écriture, shell, web, Context7, Playwright \\
\code{contract-guardian} & Relecteur en lecture seule : frontière, contrat, actions, confiance, tests ; rend un tableau de conformité & lecture, shell limité (\code{git diff}, \code{uv run}) \\
\bottomrule
\end{xltabular}}

Configuration MCP proposée dans \file{docs/kiro/mcp.example.json}, à copier dans \file{.kiro/settings/mcp.json} (l'agent n'a pas le droit d'écrire dans ce dossier) :
\begin{itemize}[leftmargin=1.2em]
  \item \textbf{Context7} : documentation à jour de \code{bpy}, React Flow, Chrome Extensions, FastAPI.
  \item \textbf{Playwright MCP} : piloter Chrome avec l'extension et le tableau de bord pendant le développement.
  \item \textbf{wysp-dev} (désactivé jusqu'à la tâche optionnelle 13 de \file{core-orchestrator}) : pont MCP en lecture seule vers un tuteur en cours d'exécution (adaptateurs, observation, tableau, audit). Kiro peut alors inspecter l'état réel de Blender pendant qu'on développe Blender.
\end{itemize}

\subsection{Utilisation au quotidien}

\begin{itemize}[leftmargin=1.2em]
  \item Ouvrir la spec de son périmètre, lancer la prochaine tâche dont on est responsable, laisser les hooks tourner, relire le diff.
  \item Référencer une spec dans le chat avec \code{\#spec:adapter-contract} ; inclure une règle manuelle avec \code{\#demo-check}.
  \item Déléguer à un agent : « utilise l'agent \code{blender-adapter} pour la tâche 4 de \file{blender-tutor} ».
  \item Avant chaque pull request : « lance \code{contract-guardian} sur ma branche ».
  \item Noter dans \file{docs/kiro-journal.md} les moments où Kiro a été décisif (spec générée, bug trouvé par un hook, propriété qui a cassé) avec la date : c'est la matière du texte de soumission et de la vidéo.
\end{itemize}
